\documentclass[../../main.tex]{subfiles}

\graphicspath{{../../image/}{../../figure/}}

\begin{document}

% TODO-BIBLIO: 本节为文献述评，方括号内为原书参考文献编号；待 reference.bib 录入对应 bib 条目后，应改为 \cite{key}。

\section{注记}\label{section:ch03-注记}

正如历史引言所指出的，非线性双曲平衡律方程组理论的源头可追溯到 19 世纪中叶，并与气体动力学一同发展。经典专著 Courant--Friedrichs [1] 以数学语言充分综述了二战结束时该领域的状况。正是这批材料的提炼奠定了现今形式化数学理论的基础。

近年来出版的大量关于双曲守恒律方程组理论与数值分析的著作，见证了这一领域的生命力。这些书在范围、风格乃至内容上互为补充，这一事实本身即表明该领域的广度。

希望先以鸟瞰方式入门的读者，可从 M.E.~Taylor [2] 的论述、Evans [2]、Hörmander [2]、Lax [7] 的教材，以及 Lax [5]、Liu [28]、Dafermos [6,10]、Bressan [15,16] 的讲义入手。在某种程度上，上述文献中有些已显陈旧，但后两种仍属当代。事实上，Bressan [15] 列出了许多重要的公开问题，而 Bressan [16] 是配有大量信息性插图的入门书。

在理论方面，Jeffrey [2]、Rozdestvenski 与 Janenko [1]、Smoller [3] 是入门水平的早期综合性著作。较新的著作 Serre [11]、Bressan [9]、Holden 与 Risebro [2]、LeFloch [5] 把一般入门与所选方向上的深入研究结合起来。Chen 与 Wang [1] 的百科全书式论文以 Euler 方程气体动力学为跳板，广泛综述一维空间中严格双曲守恒律方程组的理论。最后，Majda [4]、Chang 与 Hsiao [3]、Li、Zhang 与 Yang [1]、Yuxi Zheng [1]、Lu [2]、Perthame [2]、Benzoni--Gavage 与 Serre [2]、Ben--Artzi 与 Falcovitz [1]、Tartar [4] 是范围较窄的专门专著。上述著作将在后续章节相关处再次引用，读者由此可了解它们各自的特色。

转向数值分析，LeVeque [1] 是入门教材，而 Godlewski 与 Raviart [1,2]、LeVeque [2] 给出更全面、更技术性的论述，并附有大量文献。其他有用的来源包括 Kröner [1]、Sod [1]、Toro [1]、Kulikovski、Pogorelov 与 Semenov [1]、Holden 与 Risebro [2]，以及 Tadmor [2] 的讲义。

另一个丰富资源是 Whitham [2] 的著作，它展示了该理论与众多不同应用之间的联系，并综述了多年来应用数学家研究波传播所发展的思想与技巧，其中许多仍有待更严格的解析发展。Zeldovich 与 Raizer [1,2]、Vincenti 与 Kruger [1] 是从物理学家视角写的优秀气体动力学入门，可用于建立直觉。

读者可通过查阅每两年举办一次的双曲问题国际会议论文集，了解过去约二十年间该领域研究活动的演变。在写作本书时已出版者，按时间顺序并以编者名义列于下：Carasso、Raviart 与 Serre [1]，Ballmann 与 Jeltsch [1]，Engquist 与 Gustafsson [1]，Donato 与 Oliveri [1]，Grove、Graham 与 Plohr [1]，Fey 与 Jeltsch [1]，Freistühler 与 Warnecke [1]，Hou 与 Tadmor [1]，Benzoni--Gavage 与 Serre [3]，Tadmor、Liu 与 Tzavaras [1]，Tatsien Li 与 Song Jiang [1]，Ancona、Bressan、Marcati 与 Marson [1]。

Serre [16] 提供了世纪之交该领域状况的一个富有洞察力的视角。

本文所用意义下的"熵"一词由 Lax [4] 引入。Greven、Keller 与 Warnecke [1] 编的书收录了关于物理与数学中各种"熵"概念的一批富有信息的文章。是否对具有旋转与 Galileo 变换对称性的方程组也存在熵，这一问题在 Sever [15,16] 中讨论。

方程组\eqref{align:ch03-Euler} 与\eqref{align:ch03-等熵Euler} 通常称为 Euler 方程组。关于其理论各方面的文献汗牛充栋，其中一部分将在后续章节介绍。就初步内容而言，读者可参阅任一流体力学教材，例如 Chorin 与 Marsden [1]。Harten [1] 以及 Harten、Lax、Levermore 与 Morokoff [1] 给出了凸熵的分类。

非线性弹性动力学的文献相对较少。Truesdell 与 Noll [1]、Antman [3] 的著作文献齐全，是很好的参考。

Cercignani [1] 的书是 Boltzmann 方程的优秀入门。关于连接气体动理论与连续介质理论这一方案的最新进展，见 Villani [1] 与 Vasseur [7] 的综述文章；另见 Berthelin 与 Vasseur [1]。

关于扩展热力学及其与动理论的关系，读者可参阅 Müller 与 Ruggeri [1] 的专著。通过截断较复杂方程组来生成较简单方程组的问题，由 Boillat 与 Ruggeri [1] 详细讨论。

关于电--热--机械学的系统发展（沿第~II 章连续介质热机械学的思路），见 Coleman 与 Dill [1]、Grot [1]。Bloom [1] 给出了涉及双曲平衡律方程组的电动力学问题的众多例子。与本书叙述最相关的是 Serre [25] 的论文；另见 Boillat [3,5]。本构方程\eqref{align:ch03-BornInfeld} 由 Born 与 Infeld [1] 提出。读者可在第~V 章找到它们的一些非凡性质及相关文献。关于磁流体力学，例如可参阅 Cabannes [1]、Jeffrey [1]、Landau 与 Lifshitz [1]、Kulikovskiy 与 Lyubimov [1] 的教材。

相对论理论是有趣双曲平衡律方程组的另一丰富源泉，本书不予展开。当流体速度可与光速相比时，Euler 方程应作修改以计入狭义相对论效应；参看 Taub [1]、Friedrichs [3]，以及 Christodoulou [1] 的书。从双曲平衡律理论视角研究这些方程已积累大量文献。为定位与广泛文献，读者可参阅 Groah、Smoller 与 Temple [2] 的近期专著，或 Smoller 与 Temple [3]、Groah、Smoller 与 Temple [1] 的综述文章；另见 Ruggeri [1,2]、Smoller 与 Temple [1,2]、Pant [1]、Jing Chen [1]。

在微分几何中，曲面等距嵌入与演化相关的问题也会出现有趣的双曲守恒律方程组，激波表现为"折点"；见\cref{section:ch18-等距浸入}%

以及 Arun 与 Prasad [1]。

一维空间中双曲守恒律方程组的更多例子将在第~VII 章给出。

\end{document}
